\section{AlphaCode 2: un modelo base más fuerte más un selector aprendido}

\subsection{Del modelo de código construido internamente a Gemini Pro}

\videofigure{figures/alphacode2-overview.jpg}{AlphaCode 2 sustituye el modelo de código original por Gemini Pro, y añade más componentes aprendidos tanto en el muestreo como en el filtrado.}{00:25:14--00:26:48}

AlphaCode 2 ya no preentrena un modelo de código desde cero, sino que aprovecha un foundation model general más fuerte, y luego hace fine-tuning de varias solver family específicas para programación competitiva. Esto delega el «conocimiento y la capacidad de codificar» al modelo base, y concentra la innovación del sistema en el muestreo diverso y el ordenamiento.

\videofigure{figures/alphacode2-contributions.jpg}{Los principales cambios de AlphaCode 2 incluyen el fine-tuning de Gemini, varias solver family y un scoring model de corrección previo al envío.}{00:26:48--00:27:42}

\begin{importantbox}{El modelo base y el buscador tienen una relación multiplicativa}
Un modelo más fuerte eleva la tasa de acierto por muestra $p$; una mejor búsqueda y selección elevan el aprovechamiento del presupuesto. Ambos factores determinan juntos la solve rate final, así que no se puede atribuir todo el avance a «un modelo más grande» o a «muestrear más».
\end{importantbox}

\subsection{El scoring model aprende «cuál candidato es más probable que sea correcto»}

\videofigure{figures/scoring-model.jpg}{Un scoring model independiente estima la corrección a partir del problema y el candidato de programa, y se usa para ordenar antes del envío.}{00:27:42--00:29:18}

Sea $y_i$ un candidato, y $s_\phi(x,y_i)$ la salida del scoring model. El sistema final no ordena por la probabilidad de generación del solver, sino que combina:

$$
\operatorname{rank}(y_i)=
\lambda s_\phi(x,y_i)+(1-\lambda)d(y_i,\mathcal S),
$$

donde $d(y_i,\mathcal S)$ representa la diversidad semántica respecto al conjunto ya seleccionado. Los candidatos con puntuación alta son más probablemente correctos, y el término de diversidad evita que todos los envíos provengan del mismo cúmulo de errores.

\begin{warningbox}{La probabilidad de generación no puede tomarse directamente como probabilidad de corrección}
Lo que el modelo prefiere suele ser código común, fluido y con patrones típicos, no código que pase todas las pruebas de casos límite. Al final del curso se señala además que los modelos de lenguaje a menudo tienen exceso de confianza en respuestas incorrectas, y que la calibración de probabilidades sigue siendo en sí misma un problema abierto.
\end{warningbox}

\videofigure{figures/alphacode2-sampling.jpg}{AlphaCode 2 agrega muestras a gran escala de varias solver family, y tras probarlas, agruparlas y puntuarlas, conserva sólo unos cuantos envíos.}{00:29:18--00:30:42}

\subsection{Lograr una solve rate más alta con menos muestras}

\videofigure{figures/solve-rate-scaling.jpg}{La solve rate de AlphaCode 2 crece con el presupuesto de muestreo, y alcanza resultados comparables o mejores que los del sistema anterior con muchas menos muestras.}{00:30:42--00:32:10}

Una conclusión importante desde el punto de vista de la ingeniería de AlphaCode 2 es que, al mejorar la proposal quality y la selector quality, no hace falta llevar siempre el número de muestras hasta un millón. Para los problemas sencillos, un presupuesto menor ya alcanza para cubrir la solución correcta; para los problemas difíciles, el muestreo adicional resulta más valioso.

\videofigure{figures/alphacode2-percentile.jpg}{En la evaluación simulada de competencia, AlphaCode 2 alcanza alrededor del percentil 85, muy por encima del AlphaCode original.}{00:37:18--00:39:02}

\videofigure{figures/alphacode2-tradeoffs.jpg}{AlphaCode 2 mejora la eficiencia de muestreo y la calidad de los candidatos, pero su costo de ejecución es alto y el método sigue dependiendo en gran medida de la verificación por ejecución del código.}{00:39:02--00:40:02}

\begin{knowledgebox}{El adaptive compute es el siguiente paso natural}
Primero, un modelo ligero estima la dificultad del problema y la confianza del candidato actual, y luego decide dinámicamente el número de muestras, la temperatura, la solver family y si se entra a una corrección en serie; esto puede bajar el costo promedio de «un millón de muestras fijas por problema» a una asignación según la necesidad.
\end{knowledgebox}

\subsection{Resumen de la sección}

AlphaCode 2, mediante un modelo base más fuerte, varias solver family y un scorer aprendido, mejora la proposal y la selection, lo que muestra que el valor del cómputo en tiempo de prueba no proviene sólo de la escala, sino también de la adaptación a la tarea y de un ordenamiento más preciso de los candidatos.
